\section{Discussion and Future Work}
\label{section:6}
Our four-level model of semantic content\,---\,and its application to evaluating the usefulness of descriptions\,---\,has practical implications for the design of accessible data representations, and theoretical implications for the relationship between visualization and natural language.

% \subsection{Implications for Visualization Design}
\subsection{Natural Language As An Interface Into Visualization}
\label{section:practicalImplications}

Divergent reader preferences for semantic content suggests that it is helpful to think of natural language\,---\,not only as an interface for constructing and exploring visualizations~\cite{srinivasan_collecting_2021,gao_datatone_2015,setlur_eviza_2016}\,---\,but also as an {interface into visualization, for \emph{understanding} the semantic content they convey}.
Under this framing, we can apply Beaudoin-Lafon's framework for evaluating interface models in terms of their descriptive, evaluative, and generative powers~\cite{beaudouin-lafon_instrumental_2000, beaudouin-lafon_designing_2004}, to bring further clarity to the practical design implications of our model.
First, our grounded theory process yielded a model with \emph{descriptive} power: it categorizes the semantic content conveyed by visualizations.
Second, our study with blind and sighted readers demonstrated our model's \emph{evaluative} power: it offered a means of comparing different levels of semantic content, thus revealing divergent preferences between these different reader groups.
Third, future work can now begin to study our model's \emph{generative} power:
% not only its technological and ethical implications for automatically generated captions (\S~\ref{section:qualResultsE}), but also its 
its implications for novel multimodal interfaces and accessible data representations.
For instance, our evaluation suggested that descriptions primarily intending to benefit sighted readers might aim to generate higher-level semantic content (\S~\ref{section:qualResultsB}), while those intending to benefit blind readers might instead focus on affording readers the option to customize and combine different content levels (\S~\ref{section:qualResultsD}), depending on their individual preferences (\S~\ref{section:qualResultsC}).
This latter path might involve automatically \ac{ARIA} tagging web-based charts to surface semantic content at Levels 1 \& 2, with human-authors conveying Level 3 content. Or, it might involve applying our model to develop and evaluate the outputs of automatic captioning systems\,---\,to probe their technological capabilities and ethical implications\,---\,in collaboration with the relevant communities (\S~\ref{section:qualResultsE}).
To facilitate this work, we have released our corpus of visualizations and labeled sentences under an open source license: 
\href{http://vis.csail.mit.edu/pubs/vis-text-model/data/}{\texttt{http://vis.csail.mit.edu/pubs/vis-text-model/data/}}.

% \subsection{Implications for Visualization Theory}
\subsection{Natural Language As Coequal With Visualization}
\label{section:theoreticalImplications}

{In closing, we turn to a discussion of our model's implications for visualization theory.}
Not only can we think of natural language as an interface into visualization (as above), but also as an interface into data itself; coequal with and complementary to visualization.
For example, some semantic content (e.g., Level 2 statistics or Level 4 explanations) may be best conveyed via language, without any reference to visual modalities~\cite{potluri_examining_2021, hearst_would_2019}, while other content (e.g., Level 3 clusters) may be uniquely suited to visual representation.
This coequal framing is not a departure from orthodox visualization theory, but rather a return to its linguistic and semiotic origins.
Indeed, at the start of his foundational \emph{Semiology of Graphics}, Jacques Bertin introduces a similar framing to formalize an idea at the heart of visualization theory: content can be conveyed not only through speaking or writing but also through the ``language'' of graphics~\cite{bertin_semiology_1983}.
While Bertin took natural language as a point of departure for formalizing a language of graphics, we have here pursued the inverse: taking visualization as occasioning a return to language.
This theoretical inversion opens avenues for future work, for which linguistic theory and semiotics are instructive~\cite{weber_towards_2019, maceachren_visual_2012, vickers_understanding_2013}.

Within the contemporary linguistic tradition, subfields like syntax, semantics, and pragmatics suggest opportunities for further analysis at each level of our model. And, since our model focuses on English sentences and canonical chart types, extensions to other languages and bespoke charts may be warranted.
Within the semiotic tradition, Christian Metz (a contemporary of Bertin's) emphasized the \emph{pluralistic} quality of graphics~\cite{campolo_signs_2020}: the semantic content conveyed by visualizations depends not only on their graphical sign-system, but also on various ``social codes'' such as education, class, expertise, and\,---\,we hasten to include\,---\,ability.
Our evaluation with blind and sighted readers (as well as work studying how charts are deployed in particular discourse contexts~\cite{hullman_visualization_2011, hullman_content_2015, lee_viral_2021, aiello_inventorizing_2020}) 
lends credence to Metz's conception of graphics as pluralistic: different readers will have different ideas about what makes visualizations meaningful (Fig.~\ref{fig:teaser}).
As a means of revealing these differences, we have here introduced a four-level model of semantic content.
We leave further elucidation of the relationship between visualization and natural language to future work.